% Generated draft: requires mathematical and source review.
\documentclass[11pt]{article}
\usepackage[margin=25mm]{geometry}
\usepackage{fontspec}
\setmainfont{Latin Modern Roman}
\usepackage{amsmath,amssymb,mathtools,unicode-math}
\setmathfont{Latin Modern Math}
\usepackage{enumitem,needspace,xurl,hyperref}
\setlength{\parindent}{0pt}
\setlength{\parskip}{5pt}
\setlength{\emergencystretch}{3em}
\newcommand{\N}{\mathbb{N}}
\newcommand{\Z}{\mathbb{Z}}
\newcommand{\Q}{\mathbb{Q}}
\newcommand{\R}{\mathbb{R}}
\newcommand{\C}{\mathbb{C}}
\newcommand{\F}{\mathbb{F}}
\newcommand{\A}{\mathbb{A}}
\newcommand{\PP}{\mathbb{P}}
\newcommand{\E}{\mathbb{E}}
\newcommand{\eps}{\varepsilon}
\newcommand{\dd}{\,\mathrm{d}}
\begin{document}
\section*{UnsolvedMath ID: 10600037}
\textbf{Source problem code:} catalogue entry
\subsection{Definitions and mathematical statement}
Fix a planar knot diagram \(D\) with two adjacent regions distinguished and omitted. A Formal Knot Theory state assigns one marker to each remaining region, with the markers incident to crossings and with each crossing marked exactly once; when a region occupies several corners, the chosen incident corner is part of the data. Thus states are the appropriate matchings between crossings and unomitted regions. Let \(\mathcal S(D)\) be this state set, and let \(C(D)\) be the knot-Floer chain module with these generators, in the coefficient and grading conventions of the intended version.

The task is to define its differential entirely by combinatorics of \(D\) and these states. Explicitly, give computable combinatorial coefficients \(n_D(s,t)\) in the required coefficient ring and a rule
\[\partial_D(s)=\sum_{t\in\mathcal S(D)}n_D(s,t)t,\]
with the appropriate coefficient extensions if the selected Floer theory uses extra variables, such that \(\partial_D^2=0\), the required Maslov grading and Alexander grading/filtration behavior holds, and the resulting chain theory has a justified comparison with the knot Floer theory represented by the source's complex. The construction must explain diagram-choice invariance in that comparison. The excerpt does not choose hat, minus or another filtered/graded variant, so no unsupported requirement that the differential preserve rather than decrease Alexander degree is imposed.

The source reports that the chain generators already admit this state description, but the differential uses higher-dimensional geometric methods. It explicitly acknowledges a combinatorial knot Floer definition via arc diagrams and still asks for the differential \emph{on the Formal Knot Theory state-generated complex}. Merely supplying another combinatorial complex with different generators does not specify that requested differential. The exact version, signs, auxiliary Heegaard data and intended chain-level identification are absent and must be fixed. The stated categorification connection is to the Alexander--Conway polynomial, with normalization depending on the chosen Floer convention.
\subsection{Short English statement}
Give a combinatorial rule for the knot Floer differential using the existing Formal Knot Theory state generators. The source already acknowledges other combinatorial Floer models and asks specifically for this state complex.
\subsection{Sources}
\begin{itemize}
\item Local UnsolvedMath registry, permanent ID \(10600037\), exact source code \texttt{catalogue entry}.
\item \emph{Unsolved Problems in Virtual Knot Theory and Combinatorial Knot Theory}, revised survey, problem-list item 37: \url{https://arxiv.org/abs/1409.2823}. Registry access 2026-07-29; source-tex extraction. The source mentions \emph{Formal Knot Theory} and arc-diagram constructions, but their embedded bibliographic citations are elided.
\item The research summary mentions spanning-tree and arc-diagram models and related work without precise references.
\end{itemize}
\end{document}
